% ─────────────────────────────────────────────────────────────
\phantomsection
\chapter{Experiments}
\label{sec:experiments}

This section describes the experimental setup, datasets, preprocessing pipeline, feature configuration, and results of the fall detection study. The theoretical foundations of the topological representation --- including persistent homology, phase-space embedding, and filtration construction --- are developed in the Methodology section, which also defines the evaluation protocols and the performance metrics used throughout. Figure~\ref{fig:pipeline} gives an overview of the complete pipeline from raw sensor measurement to fall/non-fall decision.

% ── Pipeline diagram ──────────────────────────────────────────
\begin{figure}[ht]
\centering
\begin{tikzpicture}[
  node distance = 0.55cm and 0.65cm,
  box/.style    = {rectangle, draw=black, rounded corners=3pt,
                   minimum width=1.80cm, minimum height=0.72cm,
                   align=center, font=\small, fill=white, thick},
  dbox/.style   = {rectangle, draw=black!60, rounded corners=3pt,
                   minimum width=1.80cm, minimum height=0.72cm,
                   align=center, font=\small, fill=gray!8,
                   dashed, thick},
  optbox/.style = {rectangle, draw=blue!55!black, rounded corners=5pt,
                   minimum width=4.6cm, minimum height=1.20cm,
                   align=center, font=\small, fill=blue!5, thick},
  arr/.style    = {->, >=Stealth, thick},
  oarr/.style   = {->, >=Stealth, thick, dashed, blue!55!black},
  lbl/.style    = {font=\footnotesize\bfseries, text=black!60}
]

%% ─── Row 1 : Preprocessing ──────────────────────────────────────
\node[box]                    (sensor)  {Sensor\\Data};
\node[box, right=of sensor]   (smv)     {SMV$(t)$};
\node[box, right=of smv]      (resamp)  {Resample\\50\,Hz};
\node[box, right=of resamp]   (win)     {Window\\Extract};

\draw[arr] (sensor)--(smv);
\draw[arr] (smv)--(resamp);
\draw[arr] (resamp)--(win);

%% ─── Row 2 : Feature Extraction (lambda-dependent) ──────────────
%% Placed directly below Row 1 (column-aligned: emb↔sensor, filt↔smv, …)
\node[dbox, below=1.20cm of sensor]  (emb)   {Delay\\Embed\\$(m,\tau,r)$};
\node[dbox, right=of emb]            (filt)  {Filtration\\$(c,\varepsilon)$};
\node[dbox, right=of filt]           (ph)    {Pers.\,Hom.\\$H_0,H_1$};
\node[dbox, right=of ph]             (vec)   {Vectorise\\$\phi_\lambda\!\in\!\R^{24}$};

\draw[arr] (emb)--(filt);
\draw[arr] (filt)--(ph);
\draw[arr] (ph)--(vec);

%% ─── Row 3 : Classification ──────────────────────────────────────
%% p(w) and tau_0 merged into one box to keep 4+4+4 column symmetry
\node[box, below=1.20cm of emb]  (scale)   {Std\,Scale};
\node[box, right=of scale]       (clf)     {Classifier\\(LR/SVM)};
\node[box, right=of clf]         (thrbox)  {$\hat{p}(w)\!\geq\!\tau_0$};
\node[box, right=of thrbox]      (dec)     {Decision\\Fall/ADL};

\draw[arr] (scale)--(clf);
\draw[arr] (clf)--(thrbox);
\draw[arr] (thrbox)--(dec);

%% ─── Row-to-row center connectors ────────────────────────────────
%% smv and filt are in column 2 (same x); straight vertical arrow.
\draw[arr] (smv.south) -- (filt.north);
%% filt and clf are in column 2 (same x); straight vertical arrow.
\draw[arr] (filt.south) -- (clf.north);

%% ─── Bayesian Optimisation box (left, aligned with Row 2) ───────
\node[optbox, left=0.75cm of emb]  (bayes)
  {\textbf{Bayesian Opt.\ (TPE,}\\
   $T{=}200$~\textbf{iterations)}\\[2pt]
   $\lambda^*=\argmax_{\lambda\in\Lambda}\widehat{F}_2^{\,\mathrm{CV}}(\lambda)$};

%% Horizontal arrow: bayes.east -> emb.west (same y-level)
\draw[oarr] (bayes.east) --
  node[above, font=\footnotesize\itshape]{$\lambda^*$}
  (emb.west);

%% ─── Phase background rectangles ──────────────────────────────
\begin{pgfonlayer}{background}
  \node[rounded corners=5pt, fill=cyan!7, draw=cyan!35, inner sep=5pt,
        label={[lbl]above:Preprocessing},
        fit=(sensor)(win)] {};
  \node[rounded corners=5pt, fill=orange!5, draw=orange!25, inner sep=5pt,
        label={[lbl, text=orange!65!black]above:%
          Feature Extraction~($\lambda$-dep.${}^{\dagger}$)},
        fit=(emb)(vec)] {};
  \node[rounded corners=5pt, fill=green!4, draw=green!28, inner sep=5pt,
        label={[lbl]above:Classification},
        fit=(scale)(dec)] {};
\end{pgfonlayer}

\node[font=\footnotesize, text=black!45,
      below=0.30cm of scale, anchor=north]
  {${}^{\dagger}$Theoretical basis: Methodology section.};

\end{tikzpicture}
\caption{Complete fall detection pipeline from raw inertial data to binary decision.
  \emph{Preprocessing}: SMV computation, resampling to 50~Hz, window extraction.
  \emph{Feature extraction} (dashed, $\lambda$-dependent): delay embedding,
  filtration construction, persistent homology ($H_0$, $H_1$), and
  vectorisation to $\phi_\lambda\in\R^{24}$.
  \emph{Classification}: standardisation, probabilistic classifier, and
  threshold comparison $\hat{p}(w)\geq\tau_0$ for binary decision.
  The \emph{Bayesian optimisation} box (left, blue) selects $\lambda^*$ by
  maximising cross-validated $F_2$ over training subjects via 200~TPE iterations;
  row-to-row connectors run through column~2 (SMV\,$\to$\,Filtration\,$\to$\,Classifier),
  preserving the 4\,$\times$\,3 grid symmetry.}
\label{fig:pipeline}
\end{figure}

% ─────────────────────────────────────────────────────────────
\section{Datasets}
\label{subsec:datasets}

Three publicly available benchmark datasets were used: MobiFall~\parencite{mobifall}, SisFall~\parencite{sisfall}, and FallAllD in its 40~Hz variant (hereafter FAD\_40Hz)~\parencite{fallalld}. Each dataset records inertial sensor signals from human subjects performing scripted fall events and daily activities. Together they provide diversity in sensor placement, subject age, and activity protocol, which is important for assessing whether the method generalises across acquisition conditions. Table~\ref{tab:dataset_summary} summarises the key characteristics of each dataset as used in this work.

\begin{table}[h!]
\centering
\caption{Summary of the three datasets as used in the experiments.}
\label{tab:dataset_summary}
\setlength{\tabcolsep}{4pt}\footnotesize
\begin{tabular}{lcccc}
\toprule
\textbf{Dataset} & \textbf{Subjects} & \textbf{Feature matrix}
  & \textbf{Resampling (Hz)} & \textbf{Placement} \\
\midrule
MobiFall   &  9 & $1808 \times 24$ & $87.5 \to 50$ & Trouser pocket \\
SisFall    & 24 & $7738 \times 24$ & $200  \to 50$ & Waist \\
FAD\_40Hz  & 13 & $9481 \times 24$ & $40   \to 50$ & Waist \\
\midrule
\textbf{Total} & \textbf{46} & $\mathbf{19{,}027 \times 24}$ & --- & --- \\
\bottomrule
\end{tabular}
\end{table}

\subsection{MobiFall}
\label{subsubsec:mobifall}

MobiFall v2.0 contains recordings from volunteers who carried a Samsung Galaxy S3 smartphone in their trouser pocket while performing four scripted fall types --- forward fall (hands first), forward fall onto knees, backward fall from a chair, and lateral fall from standing --- as well as nine daily activities such as walking, jogging, stair use, and car entry~\parencite{mobifall}. The nominal recording rate is 87.5~Hz. The device captures seven sensor channels: three-axis linear acceleration, three-axis angular velocity (gyroscope), three-axis magnetic field, barometric pressure, and device orientation expressed as azimuth, pitch, and roll angles.

For this study, only the accelerometer and gyroscope channels were retained. The magnetometer, barometric pressure signal, and orientation estimates were excluded for two reasons. First, these channels are absent or recorded at incompatible rates in the other two datasets, so retaining them would have prevented a methodologically consistent preprocessing pipeline across all three sources. Second, the magnetometer and barometric sensor do not directly encode the inertial dynamics of fall events; including them would increase feature dimensionality without a principled contribution to the classification objective.

The full v2.0 cohort contains participants at different stages of data collection; not all subjects have both fall and daily-activity recordings. Nine subjects (identifiers 2, 3, 4, 5, 7, 8, 9, 10, and 11) were included in this study. The remaining participants were excluded because they performed fall trials only, without the paired daily-activity recordings required for leave-one-subject-out evaluation: a held-out subject who has no daily-activity windows cannot contribute a meaningful test set to the General protocol. After windowing, the retained cohort yields a feature matrix of $1808 \times 24$.

\subsection{SisFall}
\label{subsubsec:sisfall}

SisFall is a waist-worn dataset capturing 15 fall types and 19 daily-activity types across a mixed group of young and older adults~\parencite{sisfall}. The sensor unit is mounted at the waist and houses two accelerometers and one gyroscope: the ADXL345 (dynamic range $\pm 16\,g$, 13-bit resolution, used primarily for fall detection), the MMA8451Q (dynamic range $\pm 8\,g$, 14-bit resolution), and the ITG3200 gyroscope (dynamic range $\pm 2000$\,°/s, 16-bit resolution). All channels are recorded at 200~Hz. The dataset encompasses 23 young adults and 15 older adults.

Of the two accelerometers, only the ADXL345 was used alongside the gyroscope. The ADXL345 was preferred because its higher dynamic range ($\pm 16\,g$) makes it better suited to capturing the impulsive acceleration peaks that characterise fall impacts; the MMA8451Q, with half the range, is susceptible to saturation during vigorous motions. Retaining a single accelerometer also maintains sensor-channel consistency with MobiFall and FAD\_40Hz, each of which provides only one accelerometer. Including the MMA8451Q in parallel would have introduced an asymmetry in feature construction across datasets and complicated the cross-dataset comparison.

Twenty-four subjects were retained: the 23 young adults and one older adult (subject identifier 29) who had sufficient trial counts. The remaining 14 older adult participants were excluded due to an insufficient number of recorded trials to support stable leave-one-subject-out evaluation. The original 200~Hz recording rate was downsampled to 50~Hz. The resulting feature matrix has dimensions $7738 \times 24$, making SisFall the largest dataset by subject count.

\subsection{FallAllD (FAD\_40Hz)}
\label{subsubsec:fallalld}

\paragraph{Original FallAllD dataset.}
FallAllD is a large-scale open dataset of human falls and activities of daily living collected from 15 subjects~\parencite{fallalld}. In its original release, recordings are obtained simultaneously from three sensor devices worn at the waist, wrist, and neck, each capturing four modalities: a three-axis accelerometer, a three-axis gyroscope, a three-axis magnetometer, and a barometric pressure sensor. The inertial sensors (accelerometer and gyroscope) are sampled at 238~Hz, the magnetometer at 80~Hz, and the barometer at 10~Hz. The complete activity label set comprises 135 categories, of which 44 are activities of daily living (labelled with identifiers 1--44) and 91 are distinct fall scenarios (identifiers 101--135) differing in direction, mechanism, and recovery condition. Across all subjects, sensor positions, and activities, the original release contains 26{,}420 individual recording files.

\paragraph{Selection of the 40~Hz derived variant.}
Using the original 238~Hz release directly in the present study would have imposed prohibitive computational costs. The Bayesian hyperparameter search runs 200 iterations per dataset, and each iteration requires extracting topological features --- which involve constructing and computing persistent homology of a point cloud --- from every window in the dataset under the current configuration. At 238~Hz, the number of samples per window is approximately 4.8 times larger than at 50~Hz, substantially increasing both the size of the delay-embedded point cloud and the cost of the homology computation. Furthermore, the four-modality, three-position structure of the original release creates a sensor fusion problem that falls outside the scope of the present methodology, which targets single-position inertial data.

For these reasons, the publicly available 40~Hz derived variant (FAD\_40Hz) was used instead~\parencite{fallalld}. This derived dataset has been pre-processed as follows relative to the original release: (i)~the inertial signals are downsampled from 238~Hz to 40~Hz; (ii)~only the accelerometer and gyroscope channels are retained, with the magnetometer and barometric sensors removed; (iii)~the neck-worn device data is excluded, leaving the waist and wrist positions; (iv)~the activity label set is filtered to a curated subset of fall and ADL types relevant to binary detection tasks. These modifications reduce the data volume by an order of magnitude and eliminate the multi-rate sensor-fusion problem imposed by the original release, while preserving the inertial channels and the subject-level label structure required by the present pipeline.

The 40~Hz signals were subsequently resampled to the common target rate of 50~Hz as described in Section~\ref{subsec:preprocessing}. After activity filtering and subject selection, 13 of the original 15 subjects contributed data; the two excluded subjects had an insufficient number of trials for stable leave-one-subject-out evaluation. The resulting feature matrix is $9{,}481 \times 24$, the largest in the study by sample count.

% ─────────────────────────────────────────────────────────────
\section{Data Processing and Preprocessing}
\label{subsec:preprocessing}

Before any feature computation, all sensor recordings were passed through a common preprocessing pipeline consisting of four sequential steps. This pipeline transforms raw three-axis measurements into a compact, orientation-independent scalar signal, resamples it to a common frequency, and partitions it into fixed-length windows ready for feature extraction. Figure~\ref{fig:pipeline} places this block at the start of the data path.

\textbf{Step 1 --- Signal Magnitude Vector.}
For each sensor (accelerometer and gyroscope separately), the three-axis measurement vector $\mathbf{x}(t) \in \R^3$ at time $t$ is reduced to its Euclidean norm:
\[
  \mathrm{SMV}(t) = \bigl\|\mathbf{x}(t)\bigr\|_2 = \sqrt{x_1(t)^2 + x_2(t)^2 + x_3(t)^2}.
\]
This produces one scalar time series per sensor that is invariant to device orientation on the body, eliminating the need for explicit attitude estimation.

\textbf{Step 2 --- Resampling to 50~Hz.}
The three datasets were originally recorded at 87.5~Hz (MobiFall), 200~Hz (SisFall), and 40~Hz (FAD\_40Hz). All signals were resampled to a common rate of 50~Hz, but the strategies differ by dataset because the resampling direction and the ratio of original to target rate are not uniform.

For SisFall (200~Hz $\to$ 50~Hz), the ratio is the integer 4:1, so integer-ratio decimation was applied: the SMV signal was subsampled by retaining every fourth sample. This produces a uniformly-spaced 50~Hz output without any interpolation step and without spectral distortion within the passband of interest, since the Nyquist frequency of the target rate (25~Hz) is well below the dominant frequency of fall events.

For FAD\_40Hz (40~Hz $\to$ 50~Hz), upsampling by a factor of 5:4 was required. The resampled signal was obtained by constructing a new time grid at 50~Hz and computing signal values at each new grid point by interpolating over the flanking 40~Hz samples. Because the ratio 5:4 is rational and the upsampling factor is moderate, this approach introduces no perceptible temporal smearing in the SMV magnitude signal.

For MobiFall (87.5~Hz $\to$ 50~Hz), the ratio 87.5:50 = 7:4 is rational but non-integer in both directions, so neither pure decimation nor pure upsampling is sufficient. A rational sample-rate conversion was applied, conceptually equivalent to first upsampling by a factor of 4 and then downsampling by a factor of 7. The resampling was performed on the scalar SMV signal rather than on the raw three-axis components, so the non-linear magnitude operation is not affected by the interpolation kernel.

In all three cases the resampling target of 50~Hz was selected on methodological grounds that are discussed in the Methodology section.

\textbf{Step 3 --- Window extraction.}
For each fall trial, a single window centred on the highest-magnitude point of the SMV signal was extracted. For daily-activity trials, consecutive non-overlapping windows of the same length were taken. The window length is a hyperparameter optimised during the experiment and is described in Section~\ref{subsec:hparam}.

\textbf{Step 4 --- Standard scaling.}
Each of the 24 features in the final feature vector was scaled to zero mean and unit variance. The scaling parameters were estimated from the training data only and applied to the test data, ensuring that no information from test subjects influences the normalisation.

\begin{figure}[H]
\centering
\begin{tikzpicture}[
  node distance = 0.5cm,
  io/.style       = {rectangle, draw=black!65, rounded corners=4pt,
                     minimum width=9cm, text width=8.8cm, minimum height=0.65cm,
                     align=center, font=\small, fill=gray!10, thick},
  flowstep/.style = {rectangle, draw=black, rounded corners=3pt,
                     minimum width=9cm, text width=8.8cm, minimum height=0.65cm,
                     align=center, font=\small, fill=white, thick},
  arr/.style      = {->, >=Stealth, thick}
]
\node[io] (in)
  {Raw inertial recording:\;$\mathrm{Acc}(t)\in\R^3$,\;$\mathrm{Gyr}(t)\in\R^3$};
\node[flowstep, below=0.5cm of in] (smv)
  {\textbf{Step 1\;---\;SMV:}\;$\mathrm{SMV}(t)=\|\mathbf{x}(t)\|_2$\quad(per sensor)};
\node[flowstep, below=0.5cm of smv] (res)
  {\textbf{Step 2\;---\;Resample to 50\,Hz}\;(MobiFall 87.5, SisFall 200, FAD 40\,Hz)};
\node[flowstep, below=0.5cm of res] (win)
  {\textbf{Step 3\;---\;Window Extraction}\;(length $L$;\;fall: centred on peak,\;ADL: stride $L$)};
\node[io, below=0.5cm of win] (out)
  {Labelled windows $\{(w_i,\,y_i,\,s_i)\}$};
\draw[arr] (in)--(smv);
\draw[arr] (smv)--(res);
\draw[arr] (res)--(win);
\draw[arr] (win)--(out);
\end{tikzpicture}
\caption{Preprocessing pipeline (Steps 1--3). Detailed description of each step
  is given in Section~\ref{subsec:preprocessing}.}
\label{fig:preprocessing}
\end{figure}

% ─────────────────────────────────────────────────────────────
\section{Feature Description}
\label{subsec:features}

Each window is transformed into a 24-dimensional feature vector. The features are divided into two groups: 20 persistence-based features that capture the topological shape of the motion trajectory, and 4 signal-level features that capture the gross magnitude behaviour of the accelerometer signal. The theoretical derivation of the persistence-based features is presented in the Methodology section; here only the list of features and their definitions are given.

% ── 1.3.1 Persistence-Based Features ─────────────────────────
\subsection{Persistence-Based Features}
\label{subsubsec:ph_features}

The topological representation of each window yields two sets of birth--death pairs: one for connected clusters in the trajectory ($H_0$) and one for closed loops ($H_1$). Ten scalar statistics are extracted from each set, giving 20 features in total. Table~\ref{tab:ph_features} lists these statistics and their definitions.

\begin{table}[h!]
\centering
\caption{Ten persistence-based statistics computed per homological dimension.
  Each row is computed twice: once for $H_0$ (features 1--10) and once for $H_1$
  (features 11--20).  Persistence of a topological feature is defined as its
  death scale minus its birth scale.}
\label{tab:ph_features}
\small
\begin{tabular}{clp{8cm}}
\toprule
\textbf{Index} & \textbf{Feature} & \textbf{Definition} \\
\midrule
1 / 11  & Bar count          & Number of features born and died, normalised by window length \\
2 / 12  & Entropy            & Shannon entropy of the normalised persistence values \\
3 / 13  & Maximum            & Persistence of the longest-lived feature \\
4 / 14  & Energy             & Sum of squared persistence values \\
5 / 15  & Mean               & Arithmetic mean of all persistence values \\
6 / 16  & Std deviation      & Standard deviation of the persistence values \\
7 / 17  & Median             & Median persistence \\
8 / 18  & First quartile     & 25th percentile of the persistence distribution \\
9 / 19  & Third quartile     & 75th percentile of the persistence distribution \\
10 / 20 & Third moment       & Third central moment (related to skewness) of the persistence values \\
\bottomrule
\end{tabular}
\end{table}

% ── 1.3.2 Signal-Level Features ──────────────────────────────
\subsection{Signal-Level Features}
\label{subsubsec:signal_features}

Four time-domain statistics are computed directly from the accelerometer SMV$(t)$ values within the window and appended to the feature vector, giving features 21--24. These capture the raw intensity profile of the motion event without any topological processing.

\begin{table}[h!]
\centering
\caption{Four signal-level features derived from the accelerometer magnitude
  SMV$(t)$ within each window.}
\label{tab:sig_features}
\small
\begin{tabular}{clp{8cm}}
\toprule
\textbf{Index} & \textbf{Feature} & \textbf{Definition} \\
\midrule
21 & Maximum       & Highest value of SMV$(t)$ in the window \\
22 & Std deviation & Standard deviation of SMV$(t)$ \\
23 & Mean          & Arithmetic mean of SMV$(t)$ \\
24 & Range         & Maximum minus minimum of SMV$(t)$ \\
\bottomrule
\end{tabular}
\end{table}

\begin{figure}[H]
\centering
\begin{tikzpicture}[
  node distance = 0.5cm,
  io/.style       = {rectangle, draw=black!65, rounded corners=4pt,
                     minimum width=9cm, text width=8.8cm, minimum height=0.65cm,
                     align=center, font=\small, fill=gray!10, thick},
  flowstep/.style = {rectangle, draw=black, rounded corners=3pt,
                     minimum width=9cm, text width=8.8cm, minimum height=0.65cm,
                     align=center, font=\small, fill=white, thick},
  arr/.style      = {->, >=Stealth, thick}
]
\node[io] (in)
  {Window $w\in\R^{L\times 2}$\;($\mathrm{SMV}_\mathrm{acc}$,\;$\mathrm{SMV}_\mathrm{gyr}$)};
\node[flowstep, below=0.5cm of in] (emb)
  {\textbf{Step 1\;---\;Delay Embedding}\;$(m,\tau,r)$\;$\to P_\lambda(w)\subset\R^m$};
\node[flowstep, below=0.5cm of emb] (filt)
  {\textbf{Step 2\;---\;Filtration}\;$(c,\varepsilon)$\;$\to K_c(\varepsilon)$\;
   (\texttt{Alpha} / \texttt{Rips} / \texttt{SparseRips})};
\node[flowstep, below=0.5cm of filt] (ph)
  {\textbf{Step 3\;---\;Persistent Homology}\;$\to\mathrm{PD}_0\;(H_0)$,\;$\mathrm{PD}_1\;(H_1)$};
\node[flowstep, below=0.5cm of ph] (vec)
  {\textbf{Step 4\;---\;Vectorisation}\;$\to$\;feat.\ 1--10\;($\mathrm{PD}_0$),
   \;11--20\;($\mathrm{PD}_1$),\;21--24\;(signal)};
\node[io, below=0.5cm of vec] (out)
  {Feature vector $\phi_\lambda(w)\in\R^{24}$};
\draw[arr] (in)--(emb);
\draw[arr] (emb)--(filt);
\draw[arr] (filt)--(ph);
\draw[arr] (ph)--(vec);
\draw[arr] (vec)--(out);
\end{tikzpicture}
\caption{Feature extraction pipeline (Steps 1--4). Detailed description of each step
  is given in Section~\ref{subsec:features}.}
\label{fig:features}
\end{figure}

% ─────────────────────────────────────────────────────────────
\section{Bayesian Hyperparameter Optimisation}
\label{subsec:hparam}

The topological feature computation depends on several configuration choices: the window length, the phase-space embedding parameters, and the type of topological complex. Rather than fixing these by hand or by exhaustive grid search, a Bayesian optimisation procedure was used to search for the combination that maximises fall detection performance on each dataset.

A single Optuna study was run per dataset, yielding 200 TPE iterations per dataset and 600 in total across the three datasets. Each trial proposes a TDA configuration $\lambda_t$; features are extracted under that configuration from a balanced training subset (all fall windows plus a random sample of three times as many ADL windows) and evaluated with a classifier-agnostic objective: both a logistic regression proxy and a linear SVM proxy are scored with three-fold stratified cross-validation and their $F_2$ scores are averaged. Classifier hyperparameters are kept fixed during this stage and are subsequently optimised by grid search (Section~\ref{sec:exp:classifiers}).

The search is guided by a Tree-structured Parzen Estimator (TPE). After each evaluation, the history of observations $\mathcal{H}_t=\{(\lambda_i,\hat{y}_i)\}$ is split at a quantile $q_\gamma$ into a good subset ($\hat{y}\geq q_\gamma$) and a bad subset ($\hat{y}<q_\gamma$). Kernel density estimates $\ell(\lambda)$ and $g(\lambda)$ are fitted to the configurations in the good and bad subsets respectively. The next candidate is proposed as $\lambda_{t+1}=\argmax_\lambda\,\ell(\lambda)/g(\lambda)$, the configuration most probable under the good density relative to the bad density. This ratio is proportional to the expected improvement over the current best observation (theoretical derivation in the Methodology section). The mechanism progressively focuses the search on regions of $\Lambda$ associated with high $F_2$, without evaluating all combinations exhaustively.

The averaged objective ensures that the recovered configuration $\lambda^*$ is a representation that performs well across classifier families rather than being tuned to a single model's inductive bias. A representation that simultaneously supports linear separability (beneficial for LR) and a geometrically structured metric space (beneficial for SVM with an RBF kernel) is a stronger finding than one optimised for a single classifier.

Table~\ref{tab:search_space} lists the parameters and the values considered during the search.

\begin{table}[h!]
\centering
\caption{Parameters and candidate values explored during Bayesian optimisation.
  A dash (---) indicates parameters not applicable for the Alpha complex type.}
\label{tab:search_space}
\small
\begin{tabular}{ll}
\toprule
\textbf{Parameter} & \textbf{Candidate values} \\
\midrule
Window length (\texttt{win\_sec}, seconds)
  & 1.0,\; 1.5,\; 2.0,\; 2.5,\; 3.0,\; 3.5,\; 4.0,\; 4.5,\; 5.0 \\
Embedding dimension (\texttt{dim})
  & 2,\; 3,\; 4,\; 5 \\
Embedding delay (\texttt{delay})
  & 1,\; 2,\; 3,\; 4,\; 5 \\
Stride factor (\texttt{stride\_factor})
  & 1,\; 2,\; 4 \\
Complex type (\texttt{complex\_type})
  & Alpha,\; Rips,\; SparseRips \\
Distance measure (\texttt{metric})
  & Euclidean,\; Cosine,\; Manhattan,\; Chebyshev \\
Filtration scale (\texttt{eps\_percentile})
  & 20\,\%,\; 40\,\%,\; 60\,\%,\; 80\,\% \\
\bottomrule
\end{tabular}
\end{table}

\begin{figure}[H]
\centering
\begin{tikzpicture}[
  node distance = 0.6cm,
  lbl/.style  = {font=\footnotesize\bfseries, text=blue!50!black},
  term/.style = {rectangle, draw=black!65, rounded corners=8pt,
                 minimum width=8cm, text width=7.8cm, minimum height=0.65cm,
                 align=center, font=\small, fill=gray!10, thick},
  proc/.style = {rectangle, draw=black, rounded corners=3pt,
                 minimum width=8cm, text width=7.8cm, minimum height=0.65cm,
                 align=center, font=\small, fill=white, thick},
  cond/.style = {rectangle, draw=black!60, rounded corners=3pt,
                 minimum width=8cm, text width=7.8cm, minimum height=0.65cm,
                 align=center, font=\small, fill=yellow!8, dashed, thick},
  arr/.style  = {->, >=Stealth, thick},
  larr/.style = {->, >=Stealth, thick, blue!55!black}
]
\node[term] (init)
  {Initialise:\;$T_0$ random configs\;$\to H=\{(\lambda_i,\hat{y}_i)\}$};
\node[proc, below=1.4cm of init] (eval)
  {Evaluate:\;$\hat{y}_t=\mathrm{mean}(\widehat{F}_2^{\mathrm{LR}},\widehat{F}_2^{\mathrm{SVM}})$,\;3-fold CV,\;update $H$};
\node[proc, below=0.6cm of eval] (update)
  {Update TPE:\;fit $\ell(\lambda)$ (good) and $g(\lambda)$ (bad)};
\node[proc, below=0.6cm of update] (propose)
  {Propose:\;$\lambda_{t+1}=\argmax_\lambda\,\ell(\lambda)/g(\lambda)$,\;
   $t\!\leftarrow\!t{+}1$};
\node[term, below=1.4cm of propose] (out)
  {Output:\;$\lambda^*=\argmax_t\,\hat{y}_t$};
%% flow
\draw[arr] (init)--(eval);
\draw[arr] (eval)--(update);
\draw[arr] (update)--(propose);
\draw[arr] (propose)--(out);
%% loop-back from propose to eval
\draw[larr] (propose.east)--++(1.2,0)
  |- node[right,font=\footnotesize\itshape,pos=0.25]{repeat}(eval.east);
%% loop background
\begin{pgfonlayer}{background}
  \node[rounded corners=5pt, draw=blue!30, fill=blue!3, inner sep=8pt,
        label={[lbl]above:{Optimisation loop ($T=200$ iterations per dataset)}},
        fit=(eval)(propose)] {};
\end{pgfonlayer}
\end{tikzpicture}
\caption{Bayesian hyperparameter optimisation loop (TPE, $T=200$ iterations per dataset).
  Each trial evaluates both an LR proxy and a linear SVM proxy; the objective is
  the mean of their $F_2$ scores, yielding a classifier-agnostic $\lambda^*$.
  Detailed description in Section~\ref{subsec:hparam}.}
\label{fig:bayesian}
\end{figure}

% ─────────────────────────────────────────────────────────────
\section{Classifiers and Training Setup}
\label{sec:exp:classifiers}

Two classifiers were compared throughout the experiments: logistic regression (LR) and a support vector machine (SVM). Both were configured with balanced class weights so that the minority fall class receives proportionally more influence during training, compensating for the natural class imbalance between fall and daily-activity windows.

Both classifiers operated on the same TDA feature representation extracted under the single optimal configuration $\lambda^*$. This shared representation reflects the classifier-agnostic optimisation objective: $\lambda^*$ is the configuration whose features support good performance across both LR and SVM rather than being tailored to either one individually.

Classifier hyperparameters were selected by \texttt{GridSearchCV} with three-fold stratified cross-validation, scoring by $F_2$. For logistic regression, the grid covered regularisation strength $C \in \{0.01,\,0.1,\,1,\,10,\,100\}$ and penalty type $\{\ell_1,\,\ell_2\}$. For the SVM, the grid covered $C \in \{0.1,\,1,\,10,\,100\}$ and kernel $\{\text{linear},\,\text{rbf}\}$; probability outputs were obtained via Platt scaling. All feature vectors were standardised (zero mean, unit variance, parameters estimated from training data only) within each cross-validation fold before fitting.

After grid-search selects the best estimator, the decision threshold $\tau$ is set by inspecting the precision--recall curve on the training fold: the threshold $\tau^* = \argmax_\tau F_2(\tau)$ that maximises $F_2$ along the curve replaces the default 0.5 cut-off. This calibration requires no additional model fitting and adapts the operating point to the recall-weighted $F_2$ objective.

% ─────────────────────────────────────────────────────────────
\section{Results}
\label{subsec:results}

The evaluation protocols (Naive, General, and Personal) and the performance metrics ($\Acc$, $\Rec$, $\Prec$, $F_1$, $F_2$) are defined in full in the Methodology section. Briefly, the General protocol uses leave-one-subject-out cross-validation and is the primary evaluation scenario; the Naive protocol uses a random 70/30 split and serves as an upper-bound reference; and the Personal protocol trains and tests on each subject's own data. Results under the General and Personal protocols are averaged over 15 repetitions and reported with 95\,\% confidence intervals. $F_2$ is the primary metric throughout.

\subsection{Optimal Hyperparameter Configurations}
\label{subsubsec:optimal_params}

Table~\ref{tab:optimal_params} shows the best configuration found by the Bayesian search for each dataset, together with the corresponding cross-validated $F_2$ score. The three datasets converged on markedly different settings.

\begin{table}[h!]
\centering
\caption{Best configuration selected by Bayesian optimisation for each dataset
  (200 iterations per dataset).  A dash (---) indicates a parameter not
  applicable for the Alpha complex.}
\label{tab:optimal_params}
\small
\begin{tabular}{lccc}
\toprule
\textbf{Parameter} & \textbf{MobiFall} & \textbf{SisFall} & \textbf{FAD\_40Hz} \\
\midrule
Best $F_2$ (search objective) & 0.9746      & 0.9733   & 0.9860      \\
Window length (s)             & 5.0         & 1.0      & 1.0         \\
Complex type                  & SparseRips  & Alpha    & Rips        \\
Embedding dimension           & 5           & 5        & 3           \\
Embedding delay               & 2           & 4        & 2           \\
Stride factor                 & 2           & 4        & 1           \\
Distance measure              & Chebyshev   & ---      & Manhattan   \\
Filtration scale              & 80\,\%      & ---      & 40\,\%      \\
\bottomrule
\end{tabular}
\end{table}

MobiFall required a long window (5.0~s) and a SparseRips complex, whereas both SisFall and FAD\_40Hz were best described by a short 1.0~s window. SisFall uniquely preferred the Alpha complex; FAD\_40Hz favoured the standard Rips construction. These differences suggest that the three datasets exhibit genuinely different motion geometries rather than a single universal pattern, which is a direct instance of Hypothesis~2 in the Methodology section.
% TODO: [Cross-reference to Discussion section once written.]

\subsection{MobiFall Results}
\label{subsubsec:mobifall_results}

Table~\ref{tab:mobifall_results} reports LR and SVM performance across the three evaluation protocols on MobiFall.

\begin{table}[h!]
\centering
\caption{MobiFall: mean performance across subjects under the three evaluation
  protocols.  General and Personal entries are mean $\pm$ 95\,\% CI over
  $n = 15$ repetitions per subject; Naive entries are means over 10 repetitions.}
\label{tab:mobifall_results}
\setlength{\tabcolsep}{3pt}\footnotesize
\begin{tabular}{llccccc}
\toprule
\textbf{Classifier} & \textbf{Protocol}
  & \textbf{Acc} & \textbf{Rec} & \textbf{Pre}
  & \textbf{$F_1$} & \textbf{$F_2$} \\
\midrule
\multirow{3}{*}{LR}
 & Naive
   & 0.9722 & 0.9512 & 0.6992 & 0.8042 & 0.8857 \\
 & General
   & $0.9774 \pm 0.0083$ & $0.9574 \pm 0.0449$ & ---
   & $0.8400 \pm 0.0487$ & $0.9049 \pm 0.0393$ \\
 & Personal
   & $0.9713 \pm 0.0098$ & $0.7319 \pm 0.0309$ & ---
   & $0.7488 \pm 0.0526$ & $0.7325 \pm 0.0295$ \\
\midrule
\multirow{3}{*}{SVM}
 & Naive
   & 0.9605 & 0.9558 & 0.6177 & 0.7461 & 0.8573 \\
 & General
   & $0.9734 \pm 0.0080$ & $0.9543 \pm 0.0404$ & ---
   & $0.8176 \pm 0.0455$ & $0.8918 \pm 0.0327$ \\
 & Personal
   & $0.9635 \pm 0.0096$ & $0.5822 \pm 0.0432$ & ---
   & $0.6371 \pm 0.0611$ & $0.5981 \pm 0.0479$ \\
\bottomrule
\end{tabular}
\end{table}

Under the General protocol LR reached $F_2 = 0.9049 \pm 0.0393$ and SVM $F_2 = 0.8918 \pm 0.0327$. The Personal protocol caused a substantial drop for both classifiers, with SVM falling to $F_2 = 0.5981$. This behaviour indicates that MobiFall provides too few trials per subject to support a reliable individual model; the subject-level sample count is simply insufficient to estimate a stable decision boundary from within-subject data alone.

\subsection{SisFall Results}
\label{subsubsec:sisfall_results}

Table~\ref{tab:sisfall_results} presents SisFall performance.

\begin{table}[h!]
\centering
\caption{SisFall: mean performance across subjects.  Format as in
  Table~\ref{tab:mobifall_results}.}
\label{tab:sisfall_results}
\setlength{\tabcolsep}{3pt}\footnotesize
\begin{tabular}{llccccc}
\toprule
\textbf{Classifier} & \textbf{Protocol}
  & \textbf{Acc} & \textbf{Rec} & \textbf{Pre}
  & \textbf{$F_1$} & \textbf{$F_2$} \\
\midrule
\multirow{3}{*}{LR}
 & Naive
   & 0.9788 & 0.9821 & 0.9308 & 0.9557 & 0.9713 \\
 & General
   & $0.9776 \pm 0.0045$ & $0.9836 \pm 0.0130$ & ---
   & $0.9529 \pm 0.0100$ & $0.9709 \pm 0.0108$ \\
 & Personal
   & $0.9656 \pm 0.0056$ & $0.9496 \pm 0.0073$ & ---
   & $0.9284 \pm 0.0113$ & $0.9405 \pm 0.0080$ \\
\midrule
\multirow{3}{*}{SVM}
 & Naive
   & 0.9790 & 0.9822 & 0.9314 & 0.9561 & 0.9715 \\
 & General
   & $0.9832 \pm 0.0040$ & $0.9742 \pm 0.0138$ & ---
   & $0.9642 \pm 0.0089$ & $0.9700 \pm 0.0111$ \\
 & Personal
   & $0.9604 \pm 0.0059$ & $0.9231 \pm 0.0096$ & ---
   & $0.9162 \pm 0.0110$ & $0.9196 \pm 0.0082$ \\
\bottomrule
\end{tabular}
\end{table}

SisFall yielded the highest General $F_2$ of the three datasets: LR achieved $0.9709 \pm 0.0108$ and SVM $0.9700 \pm 0.0111$. The two classifiers are nearly indistinguishable on this dataset, reflecting the large and well-balanced subject pool (24 subjects, 15 fall types). In the Personal protocol performance dropped moderately but remained above 0.93 for LR ($F_2 = 0.9405$) and above 0.91 for SVM ($F_2 = 0.9196$), showing that SisFall subjects contribute enough individual data to support within-subject training --- a contrast with MobiFall that illustrates how dataset size governs the viability of subject-specific models.

% TODO: [Per-subject breakdown tables for SisFall (24 subjects) exist in
%  TEZ_DataSet_TR.tex.  Confirm whether to include a condensed per-subject
%  table here or in an appendix.]

\subsection{FAD\_40Hz Results}
\label{subsubsec:fad_results}

Table~\ref{tab:fad_results} presents FAD\_40Hz performance.

\begin{table}[h!]
\centering
\caption{FAD\_40Hz: mean performance across subjects.  Format as in
  Table~\ref{tab:mobifall_results}.}
\label{tab:fad_results}
\setlength{\tabcolsep}{3pt}\footnotesize
\begin{tabular}{llccccc}
\toprule
\textbf{Classifier} & \textbf{Protocol}
  & \textbf{Acc} & \textbf{Rec} & \textbf{Pre}
  & \textbf{$F_1$} & \textbf{$F_2$} \\
\midrule
\multirow{3}{*}{LR}
 & Naive
   & 0.9802 & 0.9918 & 0.8145 & 0.8941 & 0.9502 \\
 & General
   & $0.9826 \pm 0.0042$ & $0.9929 \pm 0.0065$ & ---
   & $0.8856 \pm 0.0653$ & $0.9429 \pm 0.0363$ \\
 & Personal
   & $0.9878 \pm 0.0032$ & $0.9235 \pm 0.0450$ & ---
   & $0.9111 \pm 0.0457$ & $0.9168 \pm 0.0454$ \\
\midrule
\multirow{3}{*}{SVM}
 & Naive
   & 0.9857 & 0.9887 & 0.8624 & 0.9209 & 0.9603 \\
 & General
   & $0.9890 \pm 0.0027$ & $0.9878 \pm 0.0109$ & ---
   & $0.9236 \pm 0.0457$ & $0.9592 \pm 0.0234$ \\
 & Personal
   & $0.9839 \pm 0.0035$ & $0.8721 \pm 0.0771$ & ---
   & $0.8778 \pm 0.0592$ & $0.8724 \pm 0.0696$ \\
\bottomrule
\end{tabular}
\end{table}

On FAD\_40Hz, SVM outperformed LR in the General protocol ($F_2 = 0.9592$ vs.\ $0.9429$), reversing the ranking seen on MobiFall and SisFall. Recall remained near perfect for both classifiers ($\Rec \geq 0.987$). The wide $F_1$ and $F_2$ confidence intervals in the General protocol reflect high variability in precision across subjects, even though recall was stable --- the model consistently detects falls but the false-alarm rate differs substantially from one subject to the next.

\subsection{Cross-Dataset Performance Summary}
\label{subsubsec:summary}

Table~\ref{tab:summary} collects the mean General $F_2$ for all datasets and classifiers in a single view.

\begin{table}[h!]
\centering
\caption{Subject-level mean $F_2$ under the General (LOSO) protocol.
  Values are mean $\pm$ 95\,\% CI.}
\label{tab:summary}
\begin{tabular}{lcc}
\toprule
\textbf{Dataset} & \textbf{LR $F_2$} & \textbf{SVM $F_2$} \\
\midrule
MobiFall   & $0.9049 \pm 0.0393$ & $0.8918 \pm 0.0327$ \\
SisFall    & $0.9709 \pm 0.0108$ & $0.9700 \pm 0.0111$ \\
FAD\_40Hz  & $0.9429 \pm 0.0363$ & $0.9592 \pm 0.0234$ \\
\bottomrule
\end{tabular}
\end{table}

SisFall ranked highest for both classifiers, followed by FAD\_40Hz and then MobiFall. The relative ordering was consistent across classifiers, but the preferred classifier changed by dataset: LR was better on MobiFall and SisFall, while SVM was better on FAD\_40Hz. This suggests that the choice of classifier is secondary to dataset geometry and size, consistent with the representation-dominance hypothesis.

% TODO: [Comparison with published baseline methods to be inserted once
%  literature survey is complete --- see proposal Section 5.4 for the
%  planned baseline families (threshold-based, statistical-feature,
%  spectral-feature).]

\subsection{Effect of Decision Threshold Optimisation}
\label{subsubsec:threshold_effect}

For each held-out subject under the General protocol the decision threshold was swept between 0.05 and 0.95 (19 values). Table~\ref{tab:threshold_gains} shows the average per-subject gain in $F_2$ obtained by choosing the subject-specific optimal threshold rather than the trained default $\hat\tau_0(s)$ (the threshold selected by the PR-curve procedure on the corresponding training fold; see Section~\ref{subsubsec:threshold} of the Methodology).

\begin{table}[h!]
\centering
\caption{Average subject-level $F_2$ gain from threshold optimisation.
  $\Delta F_2(s) = F_2(\tau^*(s)) - F_2(\hat\tau_0(s))$ where $\hat\tau_0(s)$ is the
  PR-curve trained threshold for subject $s$; values averaged over held-out subjects.}
\label{tab:threshold_gains}
\small
\begin{tabular}{lcc}
\toprule
\textbf{Dataset} & \textbf{$\Delta F_2$ (LR)} & \textbf{$\Delta F_2$ (SVM)} \\
\midrule
MobiFall   & $+0.0348$ & $+0.0474$ \\
SisFall    & $+0.0079$ & $+0.0042$ \\
FAD\_40Hz  & $+0.0250$ & $+0.0121$ \\
\bottomrule
\end{tabular}
\end{table}

The gains were largest on MobiFall --- where the default threshold was furthest from optimal --- and smallest on SisFall, where the default was already near-optimal. The gains were not uniform across subjects within any dataset. On FAD\_40Hz, subject 11 showed the largest gain observed in any experiment: $\Delta F_2 = +0.1165$ for LR, lifting that subject's score from $0.8046$ to $0.9211$ by adjusting the threshold alone, without retraining the underlying model. This illustrates that a single shared threshold cannot account for the variation in decision boundary location across individuals, and that personalised threshold calibration can meaningfully recover performance for outlier subjects.

\paragraph{Computational cost of threshold personalisation.}
The trained threshold $\hat\tau_0$ enters only the final thresholding step
and plays no role in feature extraction, classifier training, or Platt calibration.
Adjusting $\hat\tau_0$ for a held-out subject therefore requires \emph{no model retraining}:
the probability scores $\{\hat{p}_i\}$ are computed once under the trained model and
the sweep evaluates $F_2(\tau,s)$ for each $\tau\in\mathcal{T}$ by a single pass over
the precomputed score vector.  The additional computational overhead per subject is
$O\bigl(|\mathcal{T}|\cdot|W_s|\bigr)$ comparisons, where $|\mathcal{T}|=19$ and
$|W_s|$ denotes the per-subject test-window count.  This is negligible against the
Bayesian search cost, which scales as
$O\bigl(T\cdot k\cdot N_{\mathrm{tr}}\cdot C_{\mathrm{PH}}\bigr)$
where $T=200$, $k$ is the number of inner-CV folds, $N_{\mathrm{tr}}$ is the training
window count, and $C_{\mathrm{PH}}$ is the per-window persistent homology cost.
On the UHeM cluster the Bayesian search consumed approximately 75~minutes per dataset;
the full threshold sweep over all subjects in a dataset ran in under one second of wall
time---a runtime ratio exceeding $10^3$.  This asymmetry makes threshold calibration an
attractive lightweight personalisation mechanism: it is applicable at inference time
from any small per-subject calibration set, incurs no retraining cost, and requires no
access to the original training data.

\paragraph{Statistical analysis of threshold heterogeneity.}
The per-subject gains $\Delta F_2(s)$ exhibit substantial inter-subject variability within
each dataset.  On FAD\_40Hz the subject-level range extends from near-zero to
$+0.1165$ (Subject~11, LR), and on MobiFall the inter-subject standard deviation of
$\Delta F_2$ is comparable to the dataset mean gain for SVM.  A purely mathematical
argument for threshold personalisation would predict a uniform benefit: if the posterior
probability function $f_\theta\circ\phi_{\lambda^*}$ is systematically miscalibrated in
the same direction for all subjects, a single shared adjustment of $\hat\tau_0$
would lift all subjects equally.  The observed heterogeneity refutes this: subjects with
near-zero gain already sit at an optimal threshold close to $0.50$, while subjects with
large gain sit far from it---a pattern inconsistent with any single global miscalibration.
The statistical evidence is that the \emph{distribution} of $\Delta F_2(s)$ is
right-skewed with a heavy upper tail (cf.\ Subject~11), and the inter-subject standard
deviation is not small relative to the mean.  Mathematically, this implies that
$\tau^*(s) = \argmax_{\tau\in\mathcal{T}}F_2(\tau,s)$ is a genuinely subject-specific
quantity: the geometry of the decision boundary $\{f_\theta(\phi_{\lambda^*}(w))=\tau\}$
varies across individuals in a way that the shared classifier $f_\theta$ cannot
fully absorb.  This is the statistical signature of Hypothesis~4 in the proposal and
motivates the distributional reporting of $\Delta F_2$ in the Results \& Discussion
section.

% TODO: [A per-subject $\Delta F_2$ distribution figure (one boxplot panel
%  per dataset) would visualise the heterogeneity of threshold gains and
%  support the personalisation hypothesis.]


